% lattice_n_dim.tex
% Section 2 -- the general bipartite octahedral lattice in dimension n.
% At v0.1: outline-shaped.  Combinatorial definition + parity / two-
% colouring + neighbour structure to be filled in during Phase 2.

\label{sec:lattice_n_dim}

\placeholder{Section body -- write during Phase 2.  Lay out the
combinatorial definition of the bipartite octahedral lattice in
arbitrary dimension $n$; verify that the $n = 3$ specialisation
reproduces Paper~I's lattice.}

\subsection{Definition}
\label{subsec:lattice_definition}

\placeholder{Vertex set, edge set, the bipartite two-colouring, and
the neighbour structure as a function of $n$.  State the lattice as
a graph with site labels in $\mathbb{Z}^n$ (or a suitable index set)
and edges given by a fixed neighbour vector set $\{\pm e_i\}$ plus
the octahedral cross-edges.}

\subsection{Parity / Two-Colouring}
\label{subsec:lattice_parity}

\placeholder{The bipartition into ``even'' and ``odd'' sublattices;
prove that the colouring is well-defined for all $n \geq 1$ and that
edges always cross colour classes.}

\subsection{Specialisation to $n = 3$}
\label{subsec:lattice_specialisation}

\placeholder{Recover Paper~I's 3-D bipartite octahedral lattice as
the $n = 3$ instance of the general construction; cite Paper~I for
the role of this specialisation as the substrate of physics.}

\subsection{Notation conventions}
\label{subsec:lattice_notation}

\placeholder{Standing notation used throughout the rest of the paper:
$\Lambda_n$ for the lattice, $V_n^{+}$ / $V_n^{-}$ for the two colour
classes, $E_n$ for the edge set, $\epsilon$ for the lattice spacing
(used in Section~\ref{sec:continuum_limits} for the limit $\epsilon
\to 0$).}
